\documentclass[10pt,a4paper]{amsart}
\usepackage[margin=19mm]{geometry}
\usepackage{amsmath,amssymb,mathtools}
\usepackage[scr=rsfs,cal=euler]{mathalfa}
\usepackage{xfrac}
\usepackage[unicode,hidelinks,bookmarks=false]{hyperref}
\newcommand{\ie}{i.e.\ }
\newcommand{\cf}{cf.\ }
\newcommand{\resp}{respectively\ }
\newcommand{\Subsection}{\subsection}
\providecommand{\nobreakdash}{\nobreak\hbox{-}}
\setlength{\emergencystretch}{2em}
\multlinegap0pt
\usepackage{bm}
\usepackage{bbm}
\usepackage{dsfont}
\usepackage{xspace}
\usepackage{cancel}
\usepackage{mathtools}
\usepackage{amsthm}
\makeatletter
\@ifundefined{theorem}{\newtheorem{theorem}{Theorem}}{}
\@ifundefined{lemma}{\newtheorem{lemma}[theorem]{Lemma}}{}
\makeatother
\DeclareMathOperator{\cont}{\lrcorner}
\newcommand{\CC}{\mathbb{C}}
\newcommand{\cI}{\mathcal{I}}
\newcommand{\rad}{ {\rm rad}}
\title[Remarks on the second Chern class of a foliation]{Remarks on the second Chern class of a foliation}
\author{Alan Muniz}
\date{Source: arXiv:2601.12558v1 (CC BY 4.0)}
\begin{document}
\maketitle

\noindent\emph{Source and licence.} Remarks on the second Chern class of a foliation. Version arXiv:2601.12558v1 (CC BY 4.0), distributed under the Creative Commons Attribution 4.0 International licence, as recorded in the arXiv licence metadata for this exact version and shown on the abstract page https://arxiv.org/abs/2601.12558v1. Full rights notice: licenses/arxiv-2601.12558-CC-BY-4.0.txt.\par\smallskip

\noindent\textit{Editorial scope.} Counted unit: the Lemma whose source label is \texttt{None} in the article (source0.tex lines 251--253, source label \texttt{None}), delivered together with the article's own complete original proof of it (source0.tex lines 255--300, one \texttt{proof} environment of 46 lines). No mathematics is written, repaired or re-derived: statement and proof are the article's own text line for line. Required context retained uncounted: none. Every definition, display and label of those lines is retained unchanged. Omitted from this package and not redistributed: 1--250, 254--254, 301--569. Nothing inside a retained excerpt is omitted or altered: every retained body is delivered exactly as the article's lines are written, so no omission receipt of any kind accompanies this package. Citations: the retained excerpts carry no citation command at all, so no bibliography entry of the article is transcribed and no reference is redistributed. Reference adaptation: none was needed and none was made; every reference inside the delivered unit resolves inside it, so no unresolved reference or citation is produced. Printed slips kept literal: no printed word, symbol, hypothesis or number of the counted statement or of its proof was altered. Layout and class: the article's own class is \texttt{amsart} with packages including amssymb, amsmath, amstext, amsfonts, amsthm, mathrsfs, mathtools, mathdots, inputenc, xcolor, tikz-cd, showkeys, hyperref; the wrapper is the lane's pinned \texttt{amsart} editorial wrapper at 10pt on A4, which restates the theorem-like environments the retained text needs and adds the article's own macro definitions that the retained text needs (\texttt{\textbackslash CC}, \texttt{\textbackslash cI}, \texttt{\textbackslash cont}, \texttt{\textbackslash rad}), with extra wrapper packages (bm, bbm, dsfont, xspace, cancel, mathtools). The delivered numbering is the unit's own. Assets: no retained span contains a figure environment, a graphics-inclusion command, a PDF-page inclusion, a table or a TikZ picture; no image file is copied into this package, the package declares no asset dependency and no image file is redistributed with it. Licence: version arXiv:2601.12558v1 of this article is distributed under the Creative Commons Attribution 4.0 International licence, as recorded in the arXiv licence metadata for this exact version and shown on the abstract page https://arxiv.org/abs/2601.12558v1. A full rights notice is delivered at licenses/arxiv-2601.12558-CC-BY-4.0.txt. Characters: every retained excerpt is pure ASCII, so the delivered engine, XeLaTeX with its default Latin Modern text font, reports no missing character.
\begin{lemma}
Let $\omega \in H^0(\Omega_X^p(kH))$. Then the ideal sheaf of the vanishing locus of $d\omega \in H^0(\widehat{\Omega}_X^p(kH))$ is locally generated by the coefficients of $\omega$ and $d\omega$. 
\end{lemma}

\begin{proof}
%First, we note that $\widehat{\Omega}_X^p(kH)$ is a quotient of $\widehat{\Omega}_{\p{N}}^p(k) = \op{N}(-p)^{\binom{N+1}{p}}$. Then, it is enough to prove the result for polynomial differential forms. Indeed, once we have proven th

First, we note that the module of differential forms on $C_X$ are obtained from polynomial differential forms on $\CC^{N+1}$ and modding out by the ideal of $X$ and its differentials. Hence, we only need to prove the result for polynomial homogeneous forms.

Let $\omega$ be a homogeneous polynomial $p$-form of total degree $k$ in the variables $x_0, \dots, x_N$. We will restrict it to the affine chart $\{x_0 = 1\}$ and compare the ideals generated by the coefficients before and after restriction. We need to fix some notation.

Given $I = \{j_1 < \dots < j_r \} \subset \{0, \dots, N\}$, we use the multi-index notation $dx_I = dx_{j_1}\wedge \dots \wedge dx_{j_r}$. For $i\not\in I$, define the sign $dx_i \wedge dx_I = (-1)^{\sigma(i,I)} dx_{I\cup \{i\}}$. Then we write $\omega = \sum_{|I| = p} A_Idx_I$, for $A_I$ homogeneous polynomials of degree $k-p$. The equation $\rad \cont \omega = 0$ translates to: 
\begin{equation}\label{eq:euler-relation}
    \forall J \mid |J| = p-1, \quad \sum_{j\not\in  J} x_j A_{J\cup \{j\}}(-1)^{\sigma(j,J)} = 0.
\end{equation}
Moreover,
\begin{equation}
    d\omega = \sum_{|J| = p+1} \left(\sum_{i\in J} \frac{\partial A_{J\setminus \{i\} }}{\partial x_i} (-1)^{\sigma(i,J\setminus \{i\})} \right) dx_J. 
\end{equation}

Next, we restrict to the affine chart $\{x_0 = 1\}$ and compare the ideals defined by $d\omega$. Let $a_J = A_J|_{x_0=1}$. Note that restricting $\omega$ and $d\omega$ to $\{x_0 = 1\}$ eliminates every term with a $dx_0$. Let $\cI\in \CC[x_1, \dots, x_N]$ denote the ideal generated by the coefficients of $d\omega$ modulo $(x_0 -1)$, and let $\cI_0$ be the ideal generated by the coefficients of $d\omega|_{x_0=1}$. Then
\begin{align*}
    \cI_0 &= \left(\sum_{i\in J} \frac{\partial a_{J\setminus \{i\} }}{\partial x_i} (-1)^{\sigma(i,J\setminus \{i\})}  \,\Big|\, |J| = p+1 , \, 0 \not\in J  \right), \\
    \cI &= \cI_0 + \left(\frac{\partial A_{J\setminus \{0\} }}{\partial x_0}\Big|_{x_0= 1} + \sum_{i\in J \setminus \{0\}} \frac{\partial a_{J\setminus \{i\} }}{\partial x_i} (-1)^{\sigma(i,J\setminus \{i\})} \,\Big|\, 0 \in J  \right).
\end{align*}
We want to show that $\cI = \cI_0 + (a_K \mid 0\not\in K)$. Fix $K\in \{1, \dots , N\}$ such that $|K|=p$. We will show that the term in display above for $J = \{0\}\cup K$ is of the form $k\,a_K + \cI_0$. By the Euler relation and the contraction equation \eqref{eq:euler-relation}, one gets
\begin{align*}
    \frac{\partial A_{K}}{\partial x_0}\Big|_{x_0= 1} & = (k-p)a_K -\sum_{j=1}^N x_j\frac{\partial a_{K}}{\partial x_j} ,\\
    \frac{\partial a_{\{0\} \cup K  \setminus \{i\} }}{\partial x_i} &= -a_K(-1)^{\sigma(i,K \setminus \{i\})} - \sum_{j \not\in \{0\} \cup K\setminus \{i\}} x_j\frac{\partial a_{\{j\} \cup K\setminus \{i\}}}{\partial x_i}(-1)^{\sigma(j,K\setminus \{i\})}.
\end{align*}
Therefore, combining these equations, we have
% \[
% \begin{split}
%     \frac{\partial A_{K}}{\partial x_0}|_{x_0= 1} + \sum_{i\in K} \frac{\partial a_{\{0\} \cup K  \setminus \{i\} }}{\partial x_i} (-1)^{\sigma(i,\{0\} \cup K\setminus \{i\})} = \\
%     = (k-p)a_K -\sum_{j=1}^N x_j\frac{\partial a_{K}}{\partial x_j} + \\ \sum_{i\in K} \left(-a_K(-1)^{\sigma(i,K \setminus \{i\})} - \sum_{j \not\in \{0\} \cup K\setminus \{i\}} x_j\frac{\partial a_{\{j\} \cup K\setminus \{i\}}}{\partial x_i}(-1)^{\sigma(j,K\setminus \{i\})} \right) (-1)^{\sigma(i,\{0\}\cup K\setminus \{i\})} \\
%      = (k-p)a_K -\sum_{j=1}^N x_j\frac{\partial a_{K}}{\partial x_j} + \\ \sum_{i\in K} \left(a_K - \sum_{j \not\in \{0\} \cup K\setminus \{i\}} x_j\frac{\partial a_{\{j\} \cup K\setminus \{i\}}}{\partial x_i}(-1)^{\sigma(j,K\setminus \{i\})}(-1)^{\sigma(i,\{0\}\cup K\setminus \{i\})} \right) \\
%     = k\, a_K -\sum_{j=1}^N x_j\frac{\partial a_{K}}{\partial x_j} -\sum_{i\in K} \left(\sum_{j \not\in \{0\} \cup K\setminus \{i\}} x_j\frac{\partial a_{\{j\} \cup K\setminus \{i\}}}{\partial x_i}(-1)^{\sigma(j,K\setminus \{i\})}(-1)^{\sigma(i,\{0\}\cup K\setminus \{i\})} \right) \\
%     = k\, a_K -\sum_{j=1 }^N x_j\frac{\partial a_{K}}{\partial x_j} +\sum_{i\in K} \sum_{j \not\in \{0\} \cup K\setminus \{i\}} x_j\frac{\partial a_{\{j\} \cup K\setminus \{i\}}}{\partial x_i}(-1)^{\sigma(j,K\setminus \{i\})}(-1)^{\sigma(i,K\setminus \{i\})}  \\
%     = k\, a_K -\sum_{j\not\in \{0\}\cup K} x_j\frac{\partial a_{K}}{\partial x_j} -\sum_{i\in K} \sum_{j \not\in \{0\} \cup K} x_j\frac{\partial a_{\{j\} \cup K\setminus \{i\}}}{\partial x_i}(-1)^{\sigma(j,K)}(-1)^{\sigma(i,\{j\}\cup K\setminus \{i\})}  \\
%     = k\, a_K - \sum_{j \not\in \{0\} \cup K}(-1)^{\sigma(j,K)} x_j\sum_{i\in K\cup \{j\}} \frac{\partial a_{\{j\} \cup K\setminus \{i\}}}{\partial x_i}(-1)^{\sigma(i,\{j\}\cup K\setminus \{i\})} \\
%     = k\, a_K \mod \cI_0
% \end{split}
% \]
\begin{align*}
 &\frac{\partial A_{K}}{\partial x_0}\Big|_{x_0= 1} + \sum_{i\in K} \frac{\partial a_{\{0\} \cup K  \setminus \{i\} }}{\partial x_i} (-1)^{\sigma(i,\{0\} \cup K\setminus \{i\})} = \\
  &\quad     = k\, a_K - \sum_{j \not\in \{0\} \cup K}(-1)^{\sigma(j,K)} x_j\sum_{i\in K\cup \{j\}} \frac{\partial a_{\{j\} \cup K\setminus \{i\}}}{\partial x_i}(-1)^{\sigma(i,\{j\}\cup K\setminus \{i\})} \\
&\quad = k\, a_K + \cI_0,
\end{align*}
which concludes the proof.
\end{proof}

\end{document}
